\chapter{Introducción}

Desde las primeras redes de escucha acústica utilizadas en conflictos armados del siglo XX hasta los avanzados sistemas de vigilancia distribuida que actualmente supervisan infraestructuras críticas, el problema fundamental se mantiene inalterado: detectar, localizar y seguir objetos o personas en un entorno específico, con la mayor fiabilidad posible y el menor coste operativo. Lo que ha cambiado drásticamente en las últimas décadas es el contexto tecnológico en el que se aborda este problema. La miniaturización de los componentes electrónicos, la disminución del coste de los sensores y la universalización de las comunicaciones inalámbricas han democratizado el acceso a capacidades de detección que, hasta los años noventa, estaban reservadas para aplicaciones militares o industriales de alto presupuesto \cite{1,2}. Este cambio en el umbral de accesibilidad ha creado un nuevo espacio de diseño en el que investigadores, ingenieros y estudiantes pueden desarrollar sistemas funcionales de vigilancia distribuida utilizando componentes comerciales de bajo coste, generando un cuerpo de conocimiento aplicado de extraordinaria relevancia para áreas tan diversas como la robótica de servicio, la seguridad perimetral, la gestión de afluencias o la domótica avanzada \cite{3,4}.

El presente Trabajo de Fin de Grado se enmarca en esta tendencia y propone el diseño, implementación y validación de un sistema de vigilancia distribuida basado en una red de sensores de barrido angular cooperativos, capaz de detectar objetos en un espacio compartido y estimar su posición mediante la combinación de las medidas de múltiples nodos. La memoria que sigue documenta el proceso completo, desde la revisión del contexto tecnológico y la definición de objetivos hasta la validación experimental del prototipo construido.

\section{Historia y evolución de los sistemas de detección y vigilancia}

El fundamento físico que sustenta la tecnología radar, es decir, la emisión de una señal, su propagación a través del espacio, la reflexión en un obstáculo y el análisis del eco resultante para deducir la posición del objeto, fue descrito teóricamente por primera vez por Heinrich Hertz en 1888, cuando demostró de manera experimental que las ondas electromagnéticas podían reflejarse en superficies metálicas de forma similar a como lo hace la luz visible \cite{5}. No obstante, la conversión de este principio en un sistema de detección práctico demandó décadas de desarrollo, impulsadas en gran medida por las necesidades de defensa aérea que se volvieron urgentes en los días previos a la Segunda Guerra Mundial. La cadena CHAIN HOME, implementada por el Reino Unido a partir de 1937, representó la primera red operativa de detección temprana de aeronaves mediante radar y fue crucial en el desarrollo de la Batalla de Inglaterra \cite{5,6}. La experiencia en el campo de batalla aceleró el avance técnico de manera notable: en menos de diez años se pasó de sistemas rudimentarios de onda continua a radares de impulsos, de antenas fijas a sistemas de exploración angular, y de la detección cualitativa de presencia a la estimación cuantitativa de distancia, velocidad y rumbo.

La posguerra dio lugar a una diversificación gradual de las aplicaciones civiles. El radar meteorológico, la gestión del tráfico aéreo y la radionavegación marítima se transformaron en infraestructuras esenciales de la civilización contemporánea, y su evolución propició innovaciones conceptuales de gran relevancia: el efecto Doppler aplicado a la identificación de objetivos en movimiento \cite{6}, el radar de apertura sintética (SAR) para la captura de imágenes de alta resolución desde plataformas aéreas o espaciales \cite{7}, y el radar FMCW (\textit{Frequency Modulated Continuous Wave}) para la medición de distancias cortas y medias en entornos terrestres \cite{8}. Cada una de estas tecnologías respondía a necesidades específicas de aplicación, pero todas compartían una arquitectura conceptual común: un emisor, un receptor y un algoritmo de procesamiento que extrae información geométrica del tiempo de vuelo o de la variación de frecuencia de la señal reflejada.

\subsection{Historia en proyectos de vigilancia distribuida}

La implementación del principio radar en la vigilancia terrestre de áreas siguió un trayecto similar al de las grandes instalaciones militares, caracterizado por la búsqueda de sistemas más compactos, económicos y de fácil despliegue. Las primeras redes de sensores para la vigilancia terrestre, que se desarrollaron durante la Guerra de Vietnam mediante dispositivos sísmicos y acústicos dispersos desde aeronaves, anticiparon conceptualmente lo que décadas después se conocería como redes de sensores inalámbricos \cite{12,13}: un conjunto de nodos autónomos, distribuidos en el espacio, que recogen información local y la envían a un punto de agregación para su procesamiento conjunto. La diferencia entre esos sistemas y los modernos era principalmente tecnológica; la electrónica de los años sesenta no permitía el procesamiento local ni las comunicaciones bidireccionales de bajo consumo, pero la arquitectura distribuida, la diversidad espacial y la fusión de datos como estrategia para superar las limitaciones de un único sensor ya estaban presentes como principios de diseño \cite{13}.

La revolución de los sistemas microelectrónicos en la década de 1990, junto con la posterior aparición de las comunicaciones inalámbricas de corto alcance estandarizadas (IEEE 802.15.4, Bluetooth, Wi-Fi), facilitó el desarrollo de las redes de sensores inalámbricos como una disciplina académica e industrial claramente definida. El trabajo pionero de Akyildiz et al. \cite{1}, publicado en 2002, organizó los fundamentos teóricos del área y estableció los modelos de referencia que han orientado la investigación en las dos décadas siguientes. En el ámbito de la vigilancia, la disponibilidad de sensores de bajo costo como los ultrasónicos, infrarrojos o de imagen, combinada con la capacidad de comunicación inalámbrica de los nodos, abrió la posibilidad de soluciones que, por primera vez, eran económicamente viables para aplicaciones no militares: vigilancia de perímetros industriales, control de afluencia en espacios públicos, detección de intrusiones en entornos domésticos \cite{3,4,18}.

\section{Estado actual y relevancia}

El contexto tecnológico actual muestra una convergencia de elementos que favorece el avance de sistemas de vigilancia distribuidos que utilizan sensores de bajo costo. En lo que respecta a los sensores de proximidad, los dispositivos que miden la distancia mediante el tiempo de vuelo, tanto ultrasónicos como ópticos, han alcanzado un grado de madurez y accesibilidad que los hace viables para proyectos de pequeña envergadura sin sacrificar las capacidades funcionales \cite{9,10}. La precisión, el alcance y la robustez de estos sensores han experimentado un progreso constante en los últimos años, ampliando el rango de aplicaciones en las que pueden ser utilizados con garantías de fiabilidad \cite{11}.

En el ámbito de la computación embebida, la llegada de plataformas de desarrollo de bajo costo con conectividad inalámbrica integrada ha eliminado la necesidad de hardware especializado para la comunicación, lo que ha reducido drásticamente tanto el costo como la complejidad en el diseño de nodos autónomos \cite{2,14}. Este cambio ha democratizado el acceso a arquitecturas distribuidas que anteriormente requerían conocimientos y recursos típicos del ámbito profesional, permitiendo que proyectos académicos aborden problemas de coordinación multi-nodo con el mismo rigor conceptual que los sistemas industriales, aunque a una escala de laboratorio.

Desde la óptica de las aplicaciones, la necesidad de sistemas de vigilancia autónomos, no invasivos y de escaso mantenimiento está en constante aumento en diversos sectores: la logística de almacenes automatizados, los entornos de atención a personas mayores, la gestión inteligente de edificios y la robótica de servicio crean escenarios en los que la detección y localización de presencia sin infraestructura óptica, y, por ende, sin las implicaciones de privacidad vinculadas a las cámaras, resulta especialmente valiosa \cite{3,4}. En este marco, los sistemas que se fundamentan en sensores acústicos o de radiofrecuencia están ocupando un nicho en expansión, favorecido por su resistencia a las condiciones de iluminación, su bajo consumo relativo y su facilidad de integración en entornos preexistentes \cite{19,20}.

\section{Motivación del proyecto}

La motivación para llevar a cabo este trabajo proviene de la identificación de una limitación esencial en los sistemas de vigilancia de nodo único: un solo sensor de barrido angular, por muy preciso que sea, proporciona únicamente una medida de distancia desde una perspectiva específica, sin la capacidad de determinar la posición absoluta del objeto detectado ni de resolver las ambigüedades geométricas que surgen cuando varios objetos se encuentran a distancias similares en diferentes direcciones. La disposición de múltiples nodos en el espacio, con posiciones y orientaciones conocidas, transforma cualitativamente las capacidades del sistema: la información de distancia de cada nodo puede integrarse para estimar la posición del objeto mediante técnicas de trilateración, y la diversidad espacial de las mediciones disminuye la probabilidad de ocultación y amplía la cobertura total del área vigilada \cite{19,21}.

No obstante, la operación conjunta de varios sensores acústicos en un entorno compartido plantea un desafío que no tiene paralelo en los sistemas de nodo único: la interferencia entre sensores. Cuando dos nodos emiten pulsos acústicos de manera simultánea o superpuesta en el tiempo, el eco que recibe cada uno de ellos puede incluir contribuciones del pulso del otro, lo que genera lecturas incorrectas que afectan la fiabilidad del sistema en su totalidad \cite{9,11}. Este inconveniente, que en la literatura se aborda comúnmente a través de la separación temporal de las emisiones \cite{20,21}, demanda un mecanismo de coordinación explícito que asegure que los nodos con potencial de interferencia mutua no operen al mismo tiempo. El diseño de dicho mecanismo de coordinación, junto con la validación de su efectividad en un prototipo real, constituye el núcleo técnico de este estudio.

La selección de este problema como tema de un Trabajo de Fin de Grado se justifica, además, por su naturaleza como ejercicio de síntesis de competencias: el sistema resultante combina conocimientos de electrónica analógica y digital, sistemas operativos de tiempo real, protocolos de comunicación en red, algoritmos de estimación y validación experimental, abarcando de manera natural el espectro de materias del Grado en Ingeniería Informática \cite{22}.

\section{Enfoque y metodología}

El proyecto ha adoptado una metodología de desarrollo iterativo enfocada en la integración progresiva de subsistemas. A partir de un análisis de requisitos que define el alcance funcional del sistema, se ha llevado a cabo el diseño y la construcción de cada componente de hardware y software de manera incremental, verificando en cada fase la conformidad del subsistema con los requisitos correspondientes antes de su integración con el resto. Este enfoque minimiza el riesgo técnico asociado a la complejidad del sistema global al aislar los problemas en el nivel donde se originan, y permite una adaptación continua del diseño a los resultados obtenidos en las pruebas intermedias.

La validación final del sistema se ha organizado en pruebas funcionales que complementan el conjunto de requisitos identificados, siguiendo las prácticas habituales en proyectos de ingeniería de sistemas embebidos distribuidos \cite{17,18}. Cada prueba establece un escenario específico, define de antemano los criterios de éxito y documenta los resultados obtenidos de manera cuantitativa, lo que facilita la evaluación objetiva del grado de cumplimiento de los requisitos y la delimitación precisa de las capacidades del prototipo desarrollado.

\section{Estructura de la memoria}

La presente memoria se organiza en siete capítulos. El capítulo 2 establece los objetivos del proyecto y el alcance del sistema. El capítulo 3 revisa el estado del arte en sistemas de vigilancia distribuida, detección por tiempo de vuelo y coordinación de redes de sensores. El capítulo 4 formaliza los requisitos funcionales y no funcionales que condicionan el diseño. El capítulo 5 describe en detalle la arquitectura del sistema, las especificaciones hardware, el protocolo de comunicación y los mecanismos de coordinación y estimación implementados. El capítulo 6 presenta las pruebas de validación y analiza los resultados. El capítulo 7 expone las conclusiones y las líneas de trabajo futuro.
